

\documentclass[no-page-number]{nanzan-abstract}
\usepackage{float}
\usepackage{graphicx}
\usepackage{xurl}
\usepackage{balance}

\title{姿勢推定を用いた筋力トレーニング評価システムにおける処理の共通化}
\author{2023TS030 川瀬 翔大}
\professor{宮澤 元}

\begin{document}
\maketitle

\section{はじめに}
カメラ映像から人体の姿勢を推定し，関節位置や角度を利用して運動を評価するシステムが研究されている．筋力トレーニングを複数種目へ拡張する場合，種目ごとに評価対象となる関節や基準値が異なるため，どの処理を共通化し，どの部分を種目固有の処理として扱うかを設計する必要がある．
本研究では，種目ごとに異なる評価項目を設定データとして分離し，種目追加時のプログラム変更を抑えられる構成を設計する．検証する問いは，どの評価項目まで設定変更だけで扱え，どの段階から共通プログラムへの追加・変更が必要になるか，である．

\section{背景・先行研究}
田邉らは，姿勢推定を用いてサイドレイズやアームカールの運動過程を3段階に分類し，運動過程に応じたフィードバックを提示するシステムを提案している\cite{tanabe2021}．Mercadal-Baudartらは，単一カメラによる3次元姿勢推定を用いて複数の運動を定量化し，関節角度などを評価指標としている\cite{mercadal2024}．また，PatidarらはMediaPipeとOpenCVを用い，スクワット，アームカール，デッドリフト，サイドレイズなどを対象として関節角度に基づく評価を行っている\cite{patidar2025}．
これらを踏まえ，本研究では評価処理を実装する際のプログラム構成に着目し，共通処理と種目ごとの設定・個別処理をどのように分けられるかを検証する．

\section{研究概要と設計}
本研究の「複数種目」は筋力トレーニングの種目であり，まずスクワットとアームカールを具体例として扱う．スクワットでは股関節・膝・足首の3点から膝関節角度を計算し，アームカールでは肩・肘・手首の3点から肘関節角度を計算する．
3点の指定，基準値，許容範囲，比較条件を設定データに記述し，共通プログラムが設定を読み込んで特徴量計算と判定を行う．本研究の中心はJSON自体ではなく，「設定で表現する範囲」「共通処理」「個別処理」の境界を設計することである．設定だけで表現できない評価規則には個別処理を追加する．

\begin{figure}[H]
\centering
\includegraphics[width=0.95\columnwidth]{system_flow_共通化.png}
\caption{想定するシステム構成（未実装部分を含む設計案）}
\label{fig:system}
\end{figure}

\section{試作・確認}
現在はPython，OpenCV，MediaPipeを用いて試作を進めている．確認できているのは，Webカメラによる映像取得，MediaPipeによる姿勢推定，ランドマーク・骨格表示，JSONによる設定値の読み込みである．一方，関節角度の計算，設定値を用いた比較・判定，JSONによる種目切替，フィードバック表示は，現時点では動作確認済みとは扱わない．現在の試作は，共通化を検証するための準備として位置付ける．

\section{検証方法}
まずスクワットとアームカールについて，3点のランドマーク，基準値，許容範囲，比較条件を種目別設定として用意し，同一の共通プログラムで特徴量計算と判定を実行できるか確認する．その後，サイドレイズなどを追加する．
追加・変更を，(1)コード・設定とも変更不要，(2)設定のみ変更，(3)コードの追加・変更が必要，の3種類に分類する．これにより，どの評価規則まで共通処理と設定データで扱え，どこで個別処理が必要になるかを明らかにする．なお，設定した角度範囲に入ったことは，設定した評価基準に対する判定であり，「正しいフォーム」と同一視しない．

\section{考察の観点}
確認結果から，設定変更だけで扱える評価項目と，コードの追加・変更が必要な評価項目を整理する．複数関節の組合せや動作の時間的変化についても，現在の設定表現と共通処理で扱えるかを検討する．「変更量を削減できる」と主張する場合には比較対象が必要なため，比較を行わない段階では対応可能な範囲と変更箇所を記録する．

\section{今後の計画}
まず，共通化する評価項目と設定データで表現する項目を整理し，設計案を確定する．次に，3点のランドマークから角度を計算する共通処理と，JSONの基準値を用いた比較・判定処理を実装する．その後，スクワットとアームカールに同一プログラムを適用し，種目別設定だけで評価条件を切り替えられるか確認する．さらにサイドレイズなどを追加し，設定のみで対応できる範囲と個別処理が必要となる境界を検証する．

\section{おわりに}
本研究では，筋力トレーニングを題材として，複数種目を一つのシステム構成で扱うためのプログラム設計に着目する．種目ごとに異なる評価項目を設定データとして分離し，共通処理と個別処理の境界を設計した上で，どの範囲まで設定変更だけで対応できるかを検証する．現時点では映像取得，姿勢推定，ランドマーク・骨格表示，JSON読み込みまでを確認している．今後は2種目を用いた共通の特徴量計算・判定処理を実装する．

\begin{thebibliography}{3}
\bibitem{tanabe2021}
田邉英介，椋木雅之，高塚佳代子：運動過程をフィードバックする筋力トレーニング支援システム，2021年度電気・情報関係学会九州支部連合大会講演論文集，p.228，2021．
\bibitem{mercadal2024}
C. Mercadal-Baudart, C. Liu, G. Farrell, M. Boyne, J. González Escribano, A. Smolic, C. Simms:
Exercise quantification from single camera view markerless 3D pose estimation, Heliyon, Vol.10, No.6, e27596, 2024.
\bibitem{patidar2025}
K. Patidar, M. A. Zaveri:
Framework for Gym Posture Estimation Using MediaPipe and OpenCV,
2025 IEEE International Conference on Recent Advances in Computing and Systems (REACS),
DOI: 10.1109/REACS67479.2025.11413297, 2025.
\end{thebibliography}
\end{document}
